% Options for packages loaded elsewhere
\PassOptionsToPackage{unicode}{hyperref}
\PassOptionsToPackage{hyphens}{url}
%
\documentclass[
  11pt,
]{article}
\usepackage{amsmath,amssymb}
\usepackage{setspace}
\usepackage{iftex}
\ifPDFTeX
  \usepackage[T1]{fontenc}
  \usepackage[utf8]{inputenc}
  \usepackage{textcomp} % provide euro and other symbols
\else % if luatex or xetex
  \usepackage{unicode-math} % this also loads fontspec
  \defaultfontfeatures{Scale=MatchLowercase}
  \defaultfontfeatures[\rmfamily]{Ligatures=TeX,Scale=1}
\fi
\usepackage{lmodern}
\ifPDFTeX\else
  % xetex/luatex font selection
\fi
% Use upquote if available, for straight quotes in verbatim environments
\IfFileExists{upquote.sty}{\usepackage{upquote}}{}
\IfFileExists{microtype.sty}{% use microtype if available
  \usepackage[]{microtype}
  \UseMicrotypeSet[protrusion]{basicmath} % disable protrusion for tt fonts
}{}
\makeatletter
\@ifundefined{KOMAClassName}{% if non-KOMA class
  \IfFileExists{parskip.sty}{%
    \usepackage{parskip}
  }{% else
    \setlength{\parindent}{0pt}
    \setlength{\parskip}{6pt plus 2pt minus 1pt}}
}{% if KOMA class
  \KOMAoptions{parskip=half}}
\makeatother
\usepackage{xcolor}
\usepackage[margin=2.5cm]{geometry}
\usepackage{color}
\usepackage{fancyvrb}
\newcommand{\VerbBar}{|}
\newcommand{\VERB}{\Verb[commandchars=\\\{\}]}
\DefineVerbatimEnvironment{Highlighting}{Verbatim}{commandchars=\\\{\}}
% Add ',fontsize=\small' for more characters per line
\usepackage{framed}
\definecolor{shadecolor}{RGB}{248,248,248}
\newenvironment{Shaded}{\begin{snugshade}}{\end{snugshade}}
\newcommand{\AlertTok}[1]{\textcolor[rgb]{0.94,0.16,0.16}{#1}}
\newcommand{\AnnotationTok}[1]{\textcolor[rgb]{0.56,0.35,0.01}{\textbf{\textit{#1}}}}
\newcommand{\AttributeTok}[1]{\textcolor[rgb]{0.13,0.29,0.53}{#1}}
\newcommand{\BaseNTok}[1]{\textcolor[rgb]{0.00,0.00,0.81}{#1}}
\newcommand{\BuiltInTok}[1]{#1}
\newcommand{\CharTok}[1]{\textcolor[rgb]{0.31,0.60,0.02}{#1}}
\newcommand{\CommentTok}[1]{\textcolor[rgb]{0.56,0.35,0.01}{\textit{#1}}}
\newcommand{\CommentVarTok}[1]{\textcolor[rgb]{0.56,0.35,0.01}{\textbf{\textit{#1}}}}
\newcommand{\ConstantTok}[1]{\textcolor[rgb]{0.56,0.35,0.01}{#1}}
\newcommand{\ControlFlowTok}[1]{\textcolor[rgb]{0.13,0.29,0.53}{\textbf{#1}}}
\newcommand{\DataTypeTok}[1]{\textcolor[rgb]{0.13,0.29,0.53}{#1}}
\newcommand{\DecValTok}[1]{\textcolor[rgb]{0.00,0.00,0.81}{#1}}
\newcommand{\DocumentationTok}[1]{\textcolor[rgb]{0.56,0.35,0.01}{\textbf{\textit{#1}}}}
\newcommand{\ErrorTok}[1]{\textcolor[rgb]{0.64,0.00,0.00}{\textbf{#1}}}
\newcommand{\ExtensionTok}[1]{#1}
\newcommand{\FloatTok}[1]{\textcolor[rgb]{0.00,0.00,0.81}{#1}}
\newcommand{\FunctionTok}[1]{\textcolor[rgb]{0.13,0.29,0.53}{\textbf{#1}}}
\newcommand{\ImportTok}[1]{#1}
\newcommand{\InformationTok}[1]{\textcolor[rgb]{0.56,0.35,0.01}{\textbf{\textit{#1}}}}
\newcommand{\KeywordTok}[1]{\textcolor[rgb]{0.13,0.29,0.53}{\textbf{#1}}}
\newcommand{\NormalTok}[1]{#1}
\newcommand{\OperatorTok}[1]{\textcolor[rgb]{0.81,0.36,0.00}{\textbf{#1}}}
\newcommand{\OtherTok}[1]{\textcolor[rgb]{0.56,0.35,0.01}{#1}}
\newcommand{\PreprocessorTok}[1]{\textcolor[rgb]{0.56,0.35,0.01}{\textit{#1}}}
\newcommand{\RegionMarkerTok}[1]{#1}
\newcommand{\SpecialCharTok}[1]{\textcolor[rgb]{0.81,0.36,0.00}{\textbf{#1}}}
\newcommand{\SpecialStringTok}[1]{\textcolor[rgb]{0.31,0.60,0.02}{#1}}
\newcommand{\StringTok}[1]{\textcolor[rgb]{0.31,0.60,0.02}{#1}}
\newcommand{\VariableTok}[1]{\textcolor[rgb]{0.00,0.00,0.00}{#1}}
\newcommand{\VerbatimStringTok}[1]{\textcolor[rgb]{0.31,0.60,0.02}{#1}}
\newcommand{\WarningTok}[1]{\textcolor[rgb]{0.56,0.35,0.01}{\textbf{\textit{#1}}}}
\usepackage{longtable,booktabs,array}
\usepackage{calc} % for calculating minipage widths
% Correct order of tables after \paragraph or \subparagraph
\usepackage{etoolbox}
\makeatletter
\patchcmd\longtable{\par}{\if@noskipsec\mbox{}\fi\par}{}{}
\makeatother
% Allow footnotes in longtable head/foot
\IfFileExists{footnotehyper.sty}{\usepackage{footnotehyper}}{\usepackage{footnote}}
\makesavenoteenv{longtable}
\setlength{\emergencystretch}{3em} % prevent overfull lines
\providecommand{\tightlist}{%
  \setlength{\itemsep}{0pt}\setlength{\parskip}{0pt}}
\setcounter{secnumdepth}{5}
\usepackage[]{natbib}
\bibliographystyle{plainnat}
\usepackage{bookmark}
\IfFileExists{xurl.sty}{\usepackage{xurl}}{} % add URL line breaks if available
\urlstyle{same}
\hypersetup{
  pdftitle={NFert: An R package for nitrogen balance estimation and fertilizer management in precision agriculture},
  pdfkeywords={Field crops, Fertilizer management, Nitrogen, Soil
fertility, Crop nutrition, Precision agriculture, R package},
  hidelinks,
  pdfcreator={LaTeX via pandoc}}

\title{NFert: An R package for nitrogen balance estimation and
fertilizer management in precision agriculture}
\author{true \and true \and true \and true \and true \and true}
\date{2026-01-17}

\begin{document}
\maketitle
\begin{abstract}
The NFert R package provides a comprehensive toolkit for estimating
nitrogen fertilization requirements in field crops. It incorporates a
nitrogen balance model that accounts for factors such as soil
properties, crop type, previous crops, organic amendments, and
environmental conditions. The package includes functions to calculate
crop nitrogen demand, soil nitrogen supply, leaching losses,
immobilization, and contributions from various sources. Additionally,
NFert integrates with precision agriculture techniques, enabling users
to create variable rate application maps for optimizing nitrogen
fertilization. This package aims to support sustainable and efficient
nitrogen management practices in agriculture, aligning with the
principles of integrated production and contributing to reduced
environmental impact.
\end{abstract}

{
\setcounter{tocdepth}{3}
\tableofcontents
}
\setstretch{1.5}
\section{Code metadata}\label{code-metadata}

The metadata for the code repository is provided in Table
\ref{tab:metadata}.

\begin{longtable}[]{@{}
  >{\raggedright\arraybackslash}p{(\linewidth - 4\tabcolsep) * \real{0.1053}}
  >{\raggedright\arraybackslash}p{(\linewidth - 4\tabcolsep) * \real{0.7105}}
  >{\raggedright\arraybackslash}p{(\linewidth - 4\tabcolsep) * \real{0.1842}}@{}}
\caption{Code metadata for NFert package
\{\#tbl:metadata\}}\tabularnewline
\toprule\noalign{}
\begin{minipage}[b]{\linewidth}\raggedright
Nr
\end{minipage} & \begin{minipage}[b]{\linewidth}\raggedright
Code metadata description
\end{minipage} & \begin{minipage}[b]{\linewidth}\raggedright
Value
\end{minipage} \\
\midrule\noalign{}
\endfirsthead
\toprule\noalign{}
\begin{minipage}[b]{\linewidth}\raggedright
Nr
\end{minipage} & \begin{minipage}[b]{\linewidth}\raggedright
Code metadata description
\end{minipage} & \begin{minipage}[b]{\linewidth}\raggedright
Value
\end{minipage} \\
\midrule\noalign{}
\endhead
\bottomrule\noalign{}
\endlastfoot
C1 & Current code version & v1.0 \\
C2 & Permanent link to code/repository &
\url{https://github.com/mcroci/NFert} \\
C3 & Permanent link to reproducible capsule & Not available \\
C4 & Legal code license & MIT \\
C5 & Code versioning system used & git \\
C6 & Software code languages, tools and services used & R \\
C7 & Compilation requirements, operating environments and dependencies &
R environment, depends on standard R packages (e.g., raster, stats) \\
C8 & Link to developer documentation/manual &
\url{https://github.com/mcroci/NFert} \\
C9 & Support email for questions &
\href{mailto:michele.croci@unicatt.it}{\nolinkurl{michele.croci@unicatt.it}} \\
\end{longtable}

\section{Motivation and significance}\label{motivation-and-significance}

The efficient and sustainable management of nitrogen (N) fertilization
is a critical challenge in modern agriculture. Excess nitrogen
application can lead to environmental pollution, such as nitrate
leaching into groundwater and emissions of greenhouse gases (e.g.,
nitrous oxide), while insufficient nitrogen can limit crop yields and
quality, impacting food security and farmer profitability. To address
these challenges, precise and science-based methods for determining
optimal nitrogen fertilization rates are essential. Integrated
Production systems, as defined by regional guidelines such as the
``DISCIPLINARI DI PRODUZIONE INTEGRATA'' in Emilia-Romagna, provide a
framework for reducing the environmental impact of agriculture while
maintaining productivity. Accurate nitrogen management, which considers
the complex interactions between soil, plant, and environment, is a
cornerstone of such systems.

The NFert R package was developed to provide a flexible and
comprehensive tool for calculating nitrogen balance and fertilizer
requirements in field crops, aligning with the principles of integrated
production and precision agriculture. The package is based on a
well-established nitrogen balance model that considers various factors
affecting nitrogen availability and loss in agricultural systems,
drawing upon parameters and methodologies often outlined in regional
agricultural guidelines and scientific literature.

NFert addresses several key scientific and practical problems related to
nitrogen management in agricultural systems:

\begin{enumerate}
\def\labelenumi{\arabic{enumi}.}
\item
  \textbf{Quantification of Site-Specific Nitrogen Balance}: The package
  allows users to accurately estimate the nitrogen balance for a
  specific field and crop under local conditions. This moves beyond
  generalized recommendations by incorporating site-specific data on
  soil properties (texture, organic matter, C/N ratio), climate
  (rainfall patterns influencing leaching), crop characteristics (type,
  expected yield, growth stage), previous cropping history, and organic
  fertilizer applications.
\item
  \textbf{Science-Based Fertilizer Recommendation}: By providing a
  robust nitrogen balance calculation, NFert enables the determination
  of the optimal supplemental nitrogen required from fertilizers. This
  minimizes the risk of over- or under-fertilization, leading to more
  efficient nutrient use, reduced waste, lower input costs, and
  decreased environmental pollution.
\item
  \textbf{Integration of Diverse Data Sources}: NFert provides a
  framework for integrating data from various sources, including soil
  analysis, historical yield data, climate records, and potentially
  remote sensing data, into a unified nitrogen management strategy.
\item
  \textbf{Support for Precision Agriculture Implementation}: The package
  includes functionalities specifically designed to leverage spatial
  data, such as NDVI maps, for estimating variable nitrogen application
  rates across a field. This directly supports the implementation of
  Variable Rate Application (VRA) strategies, enabling targeted nutrient
  delivery based on within-field variability.
\end{enumerate}

\subsection{Contribution to scientific
discovery}\label{contribution-to-scientific-discovery}

The NFert package contributes to scientific discovery by:

\begin{itemize}
\item
  \textbf{Facilitating Mechanistic Studies}: The modular nature of the
  package allows researchers to isolate and study the impact of specific
  factors (e.g., different soil types, varying rainfall patterns, or the
  use of different organic amendments) on individual components of the
  nitrogen balance (e.g., mineralization rates, leaching losses,
  immobilization). This supports a deeper understanding of the processes
  governing nitrogen dynamics in agricultural soils.
\item
  \textbf{Enabling Scenario Analysis}: Researchers can use NFert to
  simulate and compare different nitrogen management scenarios (e.g.,
  varying fertilizer types, application timings, or rates) and evaluate
  their potential impact on nitrogen use efficiency, crop yield, and
  environmental indicators under different environmental conditions.
\item
  \textbf{Promoting Data Integration and Model Development}: NFert
  provides a foundation for integrating diverse agricultural datasets
  and can be used as a component within more complex agro-ecosystem
  models aimed at simulating long-term nutrient cycling and
  environmental impacts.
\item
  \textbf{Bridging Research and Practice}: By translating scientific
  understanding of nitrogen dynamics into a practical computational
  tool, NFert facilitates the adoption of more precise and sustainable
  nitrogen management practices by agronomists and farmers, providing a
  feedback loop between research and real-world application.
\end{itemize}

While specific research papers citing NFert are anticipated, its design
and capabilities are directly applicable to ongoing research in nutrient
management, environmental modeling of agricultural systems, and the
development of precision agriculture technologies.

\subsection{Experimental setting (how to use the
software)}\label{experimental-setting-how-to-use-the-software}

The NFert package is an R library designed for users with a working
knowledge of the R programming environment. It is typically used within
the standard R console, RStudio, or integrated into R scripts for
automated workflows. The general experimental setting involves:

\begin{enumerate}
\def\labelenumi{\arabic{enumi}.}
\tightlist
\item
  \textbf{Installation}: Users install the package from a CRAN-like
  repository or directly from the source code repository.
\end{enumerate}

\begin{Shaded}
\begin{Highlighting}[]
\CommentTok{\# Install from GitHub (once available on CRAN, use install.packages("NFert"))}
\CommentTok{\# devtools::install\_github("mcroci/NFert")}
\FunctionTok{library}\NormalTok{(NFert)}
\end{Highlighting}
\end{Shaded}

\begin{enumerate}
\def\labelenumi{\arabic{enumi}.}
\setcounter{enumi}{1}
\item
  \textbf{Data Acquisition}: Necessary input data are collected. This
  includes standard soil analysis parameters (texture, organic matter,
  pH, Ntot, SOM, CN, exchangeable K, available P), historical yield
  data, climate data (especially rainfall), details of previous crops
  and organic fertilizer applications, and for precision agriculture
  applications, spatial data like NDVI rasters. These data are typically
  prepared as R data structures (vectors, data frames, raster objects).
\item
  \textbf{Function Calls}: Users call the relevant NFert functions,
  passing the prepared data as arguments. For a full nitrogen balance,
  the \texttt{N\_balance()} function is used, while individual functions
  like \texttt{leaching\_loss()} or \texttt{soil\_fertility()} can be
  called for specific calculations. Precision agriculture functions like
  \texttt{estimate\_N\_rate\_from\_calibration\_curve()} take raster
  objects as input.
\item
  \textbf{Output Analysis}: The functions return R objects (e.g., lists,
  data frames, raster objects) containing the calculated nitrogen
  balance components or recommended fertilizer rates. Users then
  analyze, visualize, or export these outputs for decision-making or
  further processing (e.g., generating application maps in GIS
  software).
\end{enumerate}

The package documentation provides detailed information on each
function, including parameter descriptions, units, and examples, guiding
users through the process.

\subsection{Related work in
literature}\label{related-work-in-literature}

The nitrogen balance framework implemented in NFert is based on widely
accepted principles in soil fertility and nutrient management, drawing
upon models described in classical agronomy texts and contemporary
research articles. The methodologies for estimating components like
nitrogen mineralization, leaching, and immobilization are grounded in
soil science literature and often calibrated using field experimental
data. The approaches for calculating crop nitrogen demand are based on
established crop physiological principles and nutrient uptake
coefficients, similar to those found in regional fertilization
guidelines and scientific databases.

The integration of remote sensing data for variable rate application
estimation utilizes methodologies like the NDVI-based sufficiency index
approach (e.g., the Holland \& Schepers method) and calibration curve
techniques, which are standard practices in precision agriculture
research and application \citet{Holland2009} \citet{Schepers1996}.

NFert complements other computational tools in agricultural science,
including:

\begin{itemize}
\item
  \textbf{Soil process models}: More complex models simulating detailed
  soil biogeochemical processes. NFert provides a simplified, practical
  balance approach compared to these.
\item
  \textbf{Crop growth models}: Simulation models that predict crop
  development and nutrient uptake over time. NFert can provide inputs or
  be used to validate outputs of such models.
\item
  \textbf{GIS software and spatial analysis libraries in R}: Tools
  essential for processing spatial agricultural data and utilizing the
  variable rate application outputs from NFert.
\end{itemize}

The specific algorithms and parameters within NFert are influenced by
research conducted in the Emilia-Romagna region, as reflected in the
structure and content of documents like the ``DISCIPLINARI DI PRODUZIONE
INTEGRATA,'' which provides regional coefficients and guidelines for
integrated fertilization practices.

\section{Software description}\label{software-description}

The NFert package is a collection of R functions designed for
calculating nitrogen balance and estimating fertilizer requirements in
agricultural fields. It integrates soil science, agronomy, and precision
agriculture principles to provide a comprehensive toolkit for nitrogen
management. The package leverages R's capabilities for data analysis,
visualization, and modeling, making it accessible to researchers,
agronomists, and farmers.

The NFert package serves as a computational engine for implementing a
nitrogen balance approach to fertilizer management. It provides a
structured and quantitative method for assessing the nitrogen needs of a
crop and determining the optimal amount of supplemental nitrogen
required from fertilizers. By encapsulating complex agronomic and soil
science principles into accessible R functions, NFert empowers users to
move beyond generic fertilizer recommendations and adopt more precise,
site-specific practices. This precision is crucial for optimizing crop
production while minimizing the negative environmental consequences
associated with inefficient nitrogen use. The package's design
facilitates its integration into broader agricultural data analysis
workflows, including those involving spatial data and remote sensing,
thereby supporting the implementation of precision agriculture
strategies.

\subsection{Software architecture}\label{software-architecture}

The NFert package follows a modular architecture, where individual
functions handle specific calculations related to the nitrogen balance
model. This modularity allows users flexibility in performing specific
analyses or combining functions to build a complete nitrogen balance
sheet. The core components of the package include:

\begin{itemize}
\item
  \textbf{Data Input and Processing Functions}: Functions designed to
  take raw or processed agricultural and environmental data as input.
  This includes functions that interpret soil properties (e.g.,
  \texttt{tri2()} and \texttt{tri3()} for texture classification based
  on USDA methods, \texttt{calc\_soil\_group\_and\_id\_rag()} for
  determining soil group and drainage index), process rainfall data, and
  handle crop-specific parameters.
\item
  \textbf{Nitrogen Balance Component Functions}: A set of functions,
  each dedicated to calculating a specific term in the nitrogen balance
  equation. These include:

  \begin{itemize}
  \item
    \texttt{calc\_crop\_N\_demand()}: Calculates crop nitrogen
    requirement based on expected yield and crop uptake coefficients
    (e.g., from \texttt{uptake\_table}).
  \item
    \texttt{soil\_fertility()}: Estimates soil nitrogen supply (readily
    available mineral N and mineralized N from organic matter) based on
    soil properties (Ntot, SOM, CN) and crop calendar period, using
    coefficients from tables like \texttt{coefN\_readily} and
    \texttt{coefN\_mineralised}.
  \item
    \texttt{leaching\_loss()}: Calculates nitrogen losses due to
    leaching, considering factors like winter and spring rainfall, soil
    drainage index (id\_rag), and readily available nitrogen (b1),
    referencing tables like \texttt{ca.table} and \texttt{cb.table}.
  \item
    \texttt{calc\_N\_immobilization\_loss()}: Estimates nitrogen
    immobilization loss based on soil fertility (B), oxygen
    availability, and soil drainage index (id\_rag), using correction
    factors from \texttt{ca.table} and \texttt{cb.table}.
  \item
    \texttt{nitrogen\_from\_previous\_crop\_residues()}: Quantifies
    nitrogen contribution from previous crop residues based on the
    previous crop type, referencing values from \texttt{e.table}.
  \item
    \texttt{organic\_fertilization()}: Calculates nitrogen contribution
    from organic fertilizers based on the source, application frequency,
    and quantity, using factors from \texttt{f.table}.
  \item
    \texttt{natural\_contribution()}: Estimates natural nitrogen inputs
    (e.g., atmospheric deposition) based on location and crop calendar
    period, using data from \texttt{g.table} and \texttt{coef\_time}.
  \end{itemize}
\item
  \textbf{Integrated Nitrogen Balance Function}: The
  \texttt{N\_balance()} function serves as a wrapper, calling several of
  the component functions to provide a comprehensive nitrogen balance
  calculation based on a single set of input parameters. This function
  streamlines the process for users who require a full balance sheet.
\item
  \textbf{Precision Agriculture Integration Functions}: Functions
  specifically designed to work with spatial data, such as raster
  objects containing NDVI values.
  \texttt{estimate\_N\_rate\_from\_calibration\_curve()} implements
  two-point or three-point calibration to convert NDVI to N rates, while
  \texttt{estimate\_N\_rate\_from\_holland\_schepers()} applies the
  Holland \& Schepers algorithm. These functions enable the estimation
  of variable nitrogen application rates across a field.
\item
  \textbf{Helper Functions}: Additional utility functions such as
  \texttt{calculate\_N\_fertilization()} which computes the final
  nitrogen fertilization requirement from a balance result, using the
  formula: \(N_{fert} = A - B + C1 + C2 + D - E - F_{org} - G\).
\item
  \textbf{Internal Data Tables}: The package relies heavily on internal
  data tables that store parameters and coefficients derived from
  research and regional guidelines (likely those related to the
  ``DISCIPLINARI DI PRODUZIONE INTEGRATA''). These tables are accessed
  by the functions to perform the calculations and are fundamental to
  the package's operation.
\end{itemize}

The architecture is designed to be flexible, allowing users to utilize
individual components for specific analyses or the integrated
\texttt{N\_balance()} function for a complete assessment. The precision
agriculture module builds upon similar underlying principles but
operates on spatial data structures, providing a distinct but
complementary functionality.

Figure \ref{fig:architecture} illustrates the relationships between the
different functions and components of the package.

\subsection{Software functionalities}\label{software-functionalities}

The major functionalities of the NFert package provide a comprehensive
toolkit for nitrogen management in integrated production and precision
agriculture systems:

\begin{itemize}
\item
  \textbf{Comprehensive Nitrogen Balance Calculation}: Estimates the net
  nitrogen requirement of a crop by quantifying all major inputs (soil
  supply, previous crop residues, organic fertilizers, atmospheric
  deposition) and losses (leaching, immobilization) relative to the
  crop's nitrogen demand.
\item
  \textbf{Site-Specific Nutrient Management}: Enables tailored
  fertilizer recommendations based on detailed, site-specific data
  regarding soil properties, climate, and crop characteristics, moving
  away from generalized approaches.
\item
  \textbf{Soil Property Characterization}: Includes tools for
  classifying soil texture and interpreting soil analysis results
  (organic matter, total nitrogen, C/N ratio, exchangeable potassium,
  available phosphorus) to inform nitrogen balance calculations.
\item
  \textbf{Quantification of Nitrogen Losses}: Provides methods to
  estimate nitrogen losses through key pathways like leaching,
  influenced by rainfall and soil drainage, and immobilization by soil
  microorganisms.
\item
  \textbf{Integration of Organic Inputs}: Accurately accounts for the
  nitrogen contribution from various organic fertilizers and residues
  from previous crops, promoting the efficient use of organic resources
  and reducing reliance on synthetic fertilizers.
\item
  \textbf{Precision Agriculture Capabilities}: Offers functionalities to
  utilize remote sensing data, such as NDVI, for estimating spatially
  variable nitrogen application rates across a field, supporting the
  implementation of variable rate fertilization strategies.
\item
  \textbf{Support for Integrated Production Compliance}: The
  methodologies and calculations within NFert are designed to align with
  the principles and requirements of integrated production guidelines,
  assisting users in adhering to sustainable agricultural practices.
\item
  \textbf{Modular and Flexible Design}: The package's functions can be
  used independently for specific calculations or combined for a
  complete nitrogen balance assessment, offering flexibility for
  different research and application needs.
\end{itemize}

\subsection{Sample code snippets
analysis}\label{sample-code-snippets-analysis}

Let's illustrate the usage of the package with a code snippet for
calculating the nitrogen balance of a wheat crop:

\begin{Shaded}
\begin{Highlighting}[]
\CommentTok{\# Load necessary libraries}
\FunctionTok{library}\NormalTok{(NFert)}

\CommentTok{\# Define input parameters}
\NormalTok{expected\_yield\_tons\_ha }\OtherTok{\textless{}{-}} \DecValTok{6}
\NormalTok{crop }\OtherTok{\textless{}{-}} \StringTok{"Grano tenero FF (granella)"}
\NormalTok{ccp }\OtherTok{\textless{}{-}} \StringTok{"Autumn{-}winter crop \textless{}150 days"}
\NormalTok{clay }\OtherTok{\textless{}{-}} \DecValTok{10}
\NormalTok{sand }\OtherTok{\textless{}{-}} \DecValTok{40}
\NormalTok{Ntot }\OtherTok{\textless{}{-}} \FloatTok{1.2}
\NormalTok{SOM }\OtherTok{\textless{}{-}} \DecValTok{2}
\NormalTok{CN }\OtherTok{\textless{}{-}} \DecValTok{8}
\NormalTok{oxygen\_availability }\OtherTok{\textless{}{-}} \StringTok{"Normal"}
\NormalTok{winter\_rain }\OtherTok{\textless{}{-}} \DecValTok{160}
\NormalTok{start\_spring\_rain }\OtherTok{\textless{}{-}} \DecValTok{160}
\NormalTok{prev\_crop }\OtherTok{\textless{}{-}} \StringTok{"Winter cereals straw removal"}
\NormalTok{source }\OtherTok{\textless{}{-}} \StringTok{"Cattle slurry"}
\NormalTok{fertorg\_frequency }\OtherTok{\textless{}{-}} \StringTok{"every year"}
\NormalTok{forg\_quantity }\OtherTok{\textless{}{-}} \DecValTok{100}

\CommentTok{\# Calculate nitrogen balance}
\NormalTok{n\_balance\_result }\OtherTok{\textless{}{-}} \FunctionTok{N\_balance}\NormalTok{(}
  \AttributeTok{location =} \StringTok{"Plain adjacent to urbanized areas"}\NormalTok{,}
  \AttributeTok{expected\_yield\_tons\_ha =}\NormalTok{ expected\_yield\_tons\_ha,}
  \AttributeTok{crop =}\NormalTok{ crop,}
  \AttributeTok{ccp =}\NormalTok{ ccp,}
  \AttributeTok{clay =}\NormalTok{ clay,}
  \AttributeTok{sand =}\NormalTok{ sand,}
  \AttributeTok{Ntot =}\NormalTok{ Ntot,}
  \AttributeTok{SOM =}\NormalTok{ SOM,}
  \AttributeTok{CN =}\NormalTok{ CN,}
  \AttributeTok{oxygen\_availability =}\NormalTok{ oxygen\_availability,}
  \AttributeTok{winter\_rain =}\NormalTok{ winter\_rain,}
  \AttributeTok{start\_spring\_rain =}\NormalTok{ start\_spring\_rain,}
  \AttributeTok{prev\_crop =}\NormalTok{ prev\_crop,}
  \AttributeTok{source =}\NormalTok{ source,}
  \AttributeTok{fertorg\_frequency =}\NormalTok{ fertorg\_frequency,}
  \AttributeTok{forg\_quantity =}\NormalTok{ forg\_quantity}
\NormalTok{)}

\CommentTok{\# Display the nitrogen balance components}
\FunctionTok{print}\NormalTok{(n\_balance\_result)}

\CommentTok{\# Calculate the final nitrogen fertilization requirement}
\NormalTok{required\_N }\OtherTok{\textless{}{-}} \FunctionTok{calculate\_N\_fertilization}\NormalTok{(n\_balance\_result)}
\FunctionTok{cat}\NormalTok{(}\StringTok{"Required nitrogen fertilization:"}\NormalTok{, }\FunctionTok{round}\NormalTok{(required\_N, }\DecValTok{2}\NormalTok{), }\StringTok{"kg/ha}\SpecialCharTok{\textbackslash{}n}\StringTok{"}\NormalTok{)}
\end{Highlighting}
\end{Shaded}

This snippet demonstrates how users can input relevant parameters and
obtain a comprehensive nitrogen balance calculation. The
\texttt{N\_balance()} function serves as a convenient entry point,
aggregating the results from various internal calculations to provide a
complete picture of the nitrogen dynamics for the specified scenario.
The output, \texttt{n\_balance\_result}, is a data frame containing the
calculated values for each component of the nitrogen balance:

\begin{itemize}
\tightlist
\item
  \textbf{A}: Crop nitrogen demand (kg/ha)
\item
  \textbf{B}: Total soil nitrogen supply (kg/ha)
\item
  \textbf{b1}: Readily available mineral nitrogen (kg/ha)
\item
  \textbf{b2}: Mineralizable nitrogen from soil organic matter (kg/ha)
\item
  \textbf{C1}: Winter nitrogen leaching loss (kg/ha)
\item
  \textbf{C2}: Spring nitrogen leaching loss (kg/ha)
\item
  \textbf{D}: Nitrogen immobilization loss (kg/ha)
\item
  \textbf{E}: Nitrogen from previous crop residues (kg/ha)
\item
  \textbf{Forg}: Nitrogen from organic fertilizer (kg/ha)
\item
  \textbf{G}: Natural nitrogen contribution (kg/ha)
\end{itemize}

The \texttt{calculate\_N\_fertilization()} function then computes the
final nitrogen fertilization requirement using the standard formula:

\[N_{fert} = A - B + C1 + C2 + D - E - F_{org} - G\]

The clear naming of parameters and output components enhances the
interpretability of the results, allowing users to understand the
contribution of each factor to the overall nitrogen status and determine
the required fertilizer input.

\section{Illustrative examples}\label{illustrative-examples}

Beyond the comprehensive nitrogen balance calculation, NFert offers
functionalities that can be used independently or as part of a precision
agriculture workflow. A key aspect of precision agriculture is the
ability to manage inputs variably across a field based on observed
spatial variability. NFert's functions for estimating nitrogen rates
from remote sensing data are particularly relevant here.

Let's consider an example where a user wants to generate a variable rate
nitrogen application map for a corn field based on NDVI imagery.
Assuming the user has a raster file representing the NDVI values across
the field and has determined the minimum and maximum nitrogen rates
based on field trials or local recommendations (e.g., 40 kg/ha for
high-NDVI areas and 80 kg/ha for low-NDVI areas), they can use the
\texttt{estimate\_N\_rate\_from\_calibration\_curve()} function with a
two-point calibration:

\begin{Shaded}
\begin{Highlighting}[]
\CommentTok{\# Load the NFert package and the raster library}
\FunctionTok{library}\NormalTok{(NFert)}
\FunctionTok{library}\NormalTok{(raster)}

\CommentTok{\# Load the NDVI raster (replace with the actual path to your NDVI file)}
\CommentTok{\# For this example, let\textquotesingle{}s create a dummy raster for demonstration}
\FunctionTok{set.seed}\NormalTok{(}\DecValTok{123}\NormalTok{) }\CommentTok{\# for reproducibility}
\NormalTok{ndvi\_raster }\OtherTok{\textless{}{-}} \FunctionTok{raster}\NormalTok{(}\AttributeTok{nrows=}\DecValTok{10}\NormalTok{, }\AttributeTok{ncols=}\DecValTok{10}\NormalTok{, }\AttributeTok{xmn=}\DecValTok{0}\NormalTok{, }\AttributeTok{xmx=}\DecValTok{10}\NormalTok{, }\AttributeTok{ymn=}\DecValTok{0}\NormalTok{, }\AttributeTok{ymx=}\DecValTok{10}\NormalTok{)}
\NormalTok{raster}\SpecialCharTok{::}\FunctionTok{values}\NormalTok{(ndvi\_raster) }\OtherTok{\textless{}{-}} \FunctionTok{runif}\NormalTok{(raster}\SpecialCharTok{::}\FunctionTok{ncell}\NormalTok{(ndvi\_raster), }\AttributeTok{min=}\FloatTok{0.3}\NormalTok{, }\AttributeTok{max=}\FloatTok{0.9}\NormalTok{)}

\CommentTok{\# Define minimum and maximum nitrogen rates (kg/ha) }
\CommentTok{\# corresponding to max and min NDVI respectively}
\NormalTok{min\_n\_rate }\OtherTok{\textless{}{-}} \DecValTok{40} \CommentTok{\# N rate for high NDVI (healthy, less N needed)}
\NormalTok{max\_n\_rate }\OtherTok{\textless{}{-}} \DecValTok{80} \CommentTok{\# N rate for low NDVI (less healthy, more N needed)}

\CommentTok{\# Estimate nitrogen rates using two{-}point calibration}
\NormalTok{n\_rate\_map\_2pt }\OtherTok{\textless{}{-}} \FunctionTok{estimate\_N\_rate\_from\_calibration\_curve}\NormalTok{(}
  \AttributeTok{raster =}\NormalTok{ ndvi\_raster,}
  \AttributeTok{minN =}\NormalTok{ min\_n\_rate,}
  \AttributeTok{maxN =}\NormalTok{ max\_n\_rate,}
  \AttributeTok{calibration\_type =} \StringTok{"two{-}point"}\NormalTok{,}
  \AttributeTok{plot =} \ConstantTok{TRUE} \CommentTok{\# Set to TRUE to visualize the calibration curve and N rate distribution}
\NormalTok{)}

\CommentTok{\# The output \textquotesingle{}n\_rate\_map\_2pt\textquotesingle{} is a raster layer with estimated nitrogen }
\CommentTok{\# application rates in kg/ha. This raster can then be exported (e.g., }
\CommentTok{\# using raster::writeRaster) and used in precision agriculture software}
\CommentTok{\# to create a variable rate application map for a fertilizer spreader.}
\CommentTok{\# raster::writeRaster(n\_rate\_map\_2pt, "n\_rate\_map\_2pt.tif", }
\CommentTok{\#                     format="GTiff", overwrite=TRUE)}
\end{Highlighting}
\end{Shaded}

For demonstration of the Holland \& Schepers method:

\begin{Shaded}
\begin{Highlighting}[]
\CommentTok{\# Estimate nitrogen rates using Holland \& Schepers method}
\CommentTok{\# Let\textquotesingle{}s assume a base N rate of 50 kg/ha}
\NormalTok{base\_n\_rate }\OtherTok{\textless{}{-}} \DecValTok{50}

\NormalTok{hs\_result }\OtherTok{\textless{}{-}} \FunctionTok{estimate\_N\_rate\_from\_holland\_schepers}\NormalTok{(}
  \AttributeTok{ndvi\_raster =}\NormalTok{ ndvi\_raster,}
  \AttributeTok{base\_N\_rate =}\NormalTok{ base\_n\_rate,}
  \AttributeTok{plot =} \ConstantTok{TRUE} \CommentTok{\# Set to TRUE to visualize diagnostic plots}
\NormalTok{)}

\CommentTok{\# \textquotesingle{}hs\_result$dose\_raster\textquotesingle{} contains the estimated N rates based on the H\&S method.}
\CommentTok{\# \textquotesingle{}hs\_result$sufficiency\_index\_raster\textquotesingle{} contains the calculated sufficiency index.}
\CommentTok{\# These raster layers can also be exported for use in precision agriculture software.}
\CommentTok{\# raster::writeRaster(hs\_result$dose\_raster, "hs\_n\_rates.tif", }
\CommentTok{\#                     format="GTiff", overwrite=TRUE)}
\end{Highlighting}
\end{Shaded}

This example illustrates how NFert can directly support precision
agriculture workflows by translating remote sensing data into actionable
nitrogen application maps. The ability to choose between different
calibration methods (two-point, three-point, or Holland \& Schepers)
provides flexibility based on the available data and preferred approach.
The optional plotting functionality helps users visualize the
relationship between NDVI and estimated nitrogen rates and understand
the distribution of recommended rates across the field, facilitating the
interpretation and validation of the results before generating
application maps.

\section{Impact}\label{impact}

The NFert R package has the potential to significantly impact
agricultural research and practice, particularly in the area of nitrogen
management and precision agriculture.

\subsection{New research questions}\label{new-research-questions}

NFert facilitates the exploration of new research questions related to
site-specific nitrogen dynamics and optimized fertilization strategies
under varying environmental and management conditions. For example:

\begin{itemize}
\item
  How do the parameters of the nitrogen balance model (e.g.,
  mineralization rates, leaching coefficients) need to be adjusted or
  refined for different pedoclimatic zones within a region, and how can
  NFert be used to validate these adjustments?
\item
  Can NFert be integrated with seasonal climate forecasts to develop
  proactive, risk-adapted nitrogen fertilization plans that account for
  anticipated variations in rainfall and temperature?
\item
  What is the economic and environmental effectiveness of using
  NFert-generated variable rate nitrogen application maps compared to
  traditional uniform rate applications across diverse farming systems
  and landscapes?
\item
  How can NFert be used in conjunction with long-term field experiments
  to better understand the sustainability of different nitrogen
  management practices and their impact on soil health and productivity
  over decadal timescales?
\item
  Can the framework within NFert be extended or adapted to model the
  dynamics of other essential plant nutrients (e.g., phosphorus,
  potassium) in integrated production systems?
\end{itemize}

\subsection{Improvement of existing
research}\label{improvement-of-existing-research}

NFert improves the pursuit of existing research questions by providing a
standardized, transparent, and flexible computational platform for
nitrogen balance calculation. Researchers can easily implement and
compare different nitrogen management scenarios, analyze the sensitivity
of nitrogen balance components to various input parameters, and
integrate nitrogen dynamics into broader agro-ecosystem models. The
package's modularity allows researchers to isolate and study specific
processes (e.g., the impact of different organic amendments on nitrogen
mineralization) in detail while using a consistent framework. The
integration with spatial data and remote sensing tools enhances the
ability to conduct research on precision nitrogen management and
variable rate application, which are increasingly important areas of
study for optimizing resource use and minimizing environmental
footprint. NFert's basis on principles reflected in regional guidelines
also makes it a valuable tool for evaluating the effectiveness of such
guidelines and informing their future refinement.

\subsection{Changes in daily practice}\label{changes-in-daily-practice}

NFert can change the daily practice of agronomists and potentially
farmers by providing a more data-driven and site-specific approach to
nitrogen fertilization. Instead of relying on generalized
recommendations or empirical rules, users can calculate nitrogen
requirements based on their specific field conditions, leading to more
informed and precise decisions. For agronomists providing advisory
services, NFert can be a valuable tool for developing tailored
fertilization plans for their clients, potentially improving client
outcomes and strengthening the scientific basis of their
recommendations. The potential to generate variable rate application
maps can enable farmers with the necessary equipment to implement
precision fertilization, optimizing input use, potentially improving
profitability, and reducing the risk of environmental contamination.
While adoption requires access to necessary input data and familiarity
with the R environment, the package provides the computational engine
for implementing more sophisticated nitrogen management practices in
real-world agricultural settings.

\subsection{Widespread use and commercial
settings}\label{widespread-use-and-commercial-settings}

As an open-source R package, NFert is freely available, which
facilitates its potential for widespread use within the agricultural
research community globally and among early adopters of precision
agriculture technologies. Its utility is particularly high in regions
where integrated production guidelines are promoted or required. The
number of downloads from repositories like CRAN (once available) and
citations in future research papers will be key indicators of its
adoption. While there is no information provided about current
commercial use or spin-off companies, the functionalities offered by
NFert, particularly the ability to perform detailed nitrogen balance
calculations and generate variable rate application maps, are directly
relevant to commercial precision agriculture service providers,
agricultural consultants, and input suppliers. Companies offering soil
sampling, remote sensing data processing, and variable rate application
services could potentially integrate NFert into their service platforms
to enhance their analytical capabilities and provide more precise
recommendations to their farmer clients.

\section{Conclusions}\label{conclusions}

The NFert R package represents a valuable and timely contribution to the
field of nitrogen management in agriculture. By providing a
comprehensive, flexible, and transparent tool for calculating nitrogen
balance and estimating fertilizer requirements, the package empowers
researchers, agronomists, and potentially progressive farmers to adopt
more precise and sustainable nitrogen management practices. The
package's modular architecture facilitates both detailed investigation
of specific nitrogen cycle processes and integrated assessment of the
overall nitrogen balance. Furthermore, the inclusion of functionalities
for integrating remote sensing data underscores its relevance for modern
precision agriculture applications, enabling the generation of variable
rate fertilization maps. By aligning with the principles of integrated
production and drawing upon methodologies relevant to regional
guidelines, NFert supports efforts to optimize agricultural productivity
while minimizing the environmental footprint of nitrogen fertilization.
As the package is adopted and utilized by the scientific and
agricultural communities, it is expected to play an increasingly
important role in advancing research on nutrient management and
facilitating the implementation of data-driven, sustainable farming
systems globally.

\section{Acknowledgments}\label{acknowledgments}

The authors gratefully acknowledge the support of the Department of
Sustainable Crop Production at Università Cattolica del Sacro Cuore di
Piacenza. We also thank the agricultural extension services and farmers
who provided feedback during the development and testing of the package.
The development of NFert was supported by research activities focused on
integrated production and precision agriculture in the Emilia-Romagna
region.

\section{References}\label{references}

  \bibliography{references.bib}

\end{document}
